%%%%%%%%%%%%%%%%%%%%%%%%%%%%%%%%%%%%%%%%%%%%%%%%%%%%%%%%%%%%%%%%%%
%% Toto je inkludovaný LaTeXový soubor `cp_inst.tex'.           %%
%% ------------------------------------------------------------ %%
%% Verze 2022-05-23                                             %%
%% Vyvíjí & udržuje <stanislav.hledik@physics.slu.cz>           %%
%%%%%%%%%%%%%%%%%%%%%%%%%%%%%%%%%%%%%%%%%%%%%%%%%%%%%%%%%%%%%%%%%%

\chapter{Instalace}\label{install}%%%%%%%%%%%%%%%%%%%%%%%%%%%%%%%%%%%%%%%%%

\section{Získání balíčku \ipotname{}}\label{instgetpack}%%%%%%%%%%%%%%%%%%%

Domovská stránka balíčku \pkg{\ipotname} je \theshome, sekce \TeX. Můžete se
rozhodnout pro jeden z~následujících způsobů instalace: \uv{vše v~jednom} nebo
\uv{zdrojové texty a knihovny odděleně}. Ke stažení pro vás vhodné varianty
využijte odkazů v~textu níže. Kontrolní součty jsou v~souboru
\href{https://is.slu.cz/www/hle0002/vyuka/aid/ipothes/CHECKSUMS.md5}%
{\pgm{CHECKSUMS.md5}}.

\section{Instalace \uv{vše v~jednom}}\label{installinone}%%%%%%%%%%%%%%%%%%

Stáhněte si archiv
\href{https://is.slu.cz/www/hle0002/vyuka/aid/ipothes/vzorova\_prace.ipothes.zip}%
{\pgm{\thesnamecze .\ipotname .zip}} obsahující ve společném adresáři
\pgm{manual} po rozbalení kromě nezbytných \LaTeX ových
knihoven\footnote{Dokumentní třída, bibliografický styl, lokalizační soubory a
  systémové obrázky s~fiktivním zadáním a aktuálně se znaky Fyzikálního ústavu
  v~Opavě, Fi\-lo\-zo\-fic\-ko-pří\-ro\-do\-vě\-dec\-ké fakulty v~Opavě a
  Matematického ústavu v~Opavě.} také zdrojové soubory nutné pro kompilaci
tohoto hypertextového českého manuálu
\href{https://is.slu.cz/www/hle0002/vyuka/aid/ipothes/vzorova\_prace.pdf}%
{\pgm{\thesnamecze .pdf}}.\footnote{Není součástí výše uvedeného archivu, jenž
  neobsahuje výstupy ani intermediální soubory \pgm{*.aux}, \pgm{*.log},
  \pgm{*.toc}, \pgm{*.lof}, \pgm{*.lot}, \pgm{*.out}, \pgm{*.bbl},
  \pgm{*.blg}, \pgm{*.idx}, \pgm{*.ind}, \pgm{*.ilg}.} Tuto možnost nejspíše
zvolí uživatelé, kteří zamýšlejí využít \pkg{\ipotname} pouze pro jednorázové
napsání závěrečné práce a nechtějí si jím \uv{kontaminovat} svůj \TeX{}
systém. Je rovněž určena pro uživatele cloudového systému
Overleaf~\cite{Ovr:LIVE::Overleaf}.\footnote{Výhodou je také, že po dokončení
  práce a archivování se uchovají pohromadě zdrojové texty i~\LaTeX ové
  knihovny (jež budou pravděpodobně vyvíjeny bez záruky zpětné kompatibility),
  takže kompatibilita je víceméně zachována. Nevýhodou je vyšší míra chaosu ve
  vašem pracovním adresáři.}

Váš \TeX{} systém není třeba v~tomto případě nijak konfigurovat, všechny
soubory budou nalezeny. V~pracovním adresáři \pgm{manual} (klidně si jej
přejmenujte) si vytvořte nový vlastní hlavní \LaTeX ový soubor a případně
i~inkludované soubory \pgm{include/*.tex} podle ukázkového hlavního \LaTeX
ového souboru \pgm{\thesnamecze .tex}\footnote{Velmi doporučujeme, aby
  v~názvech vašich zdrojových souborů ani v~cestách k~nim nebyly znaky
  s~diakritikou a mezery. Ušetříte si tím možná nepříjemná překvapení.} nebo
pokračujte formou přejmenovávání a editace obsahu zdrojáků tohoto manuálu.

\section{Instalace \uv{zdrojové texty a knihovny odděleně}}\label{instsepar}%%

Stáhněte si archivy

\begin{enumerate}
\item \href{https://is.slu.cz/www/hle0002/vyuka/aid/ipothes/ipothes.tds.zip}%
  {\pgm{\ipotname .tds.zip}} obsahující pouze knihovní
  soubory\footnote{Dokumentní třída, bibliografický styl, dokumentační
    soubory, lokalizační soubory a systémové obrázky s~fiktivním zadáním a
    aktuálně se znaky Fyzikálního ústavu v~Opavě,
    Fi\-lo\-zo\-fic\-ko-pří\-ro\-do\-vě\-dec\-ké fakulty v~Opavě a
    Matematického ústavu v~Opavě.} uložené v~kanonické adresářové struktuře
  požadované systémem \TeX~\cite{Sch:LIVE::TDS}, a to v~adresářích
  \pgm{bibtex}, \pgm{doc} a \pgm{tex} a struktuře pod nimi;
\item
  \href{https://is.slu.cz/www/hle0002/vyuka/aid/ipothes/vzorova\_prace.zip}%
  {\pgm{\thesnamecze .zip}} obsahující v~adresáři \pgm{manual} pouze zdrojové
  soubory nutné ke kompilaci tohoto českého manuálu
  \href{https://is.slu.cz/www/hle0002/vyuka/aid/ipothes/vzorova\_prace.pdf}%
  {\pgm{\thesnamecze .pdf}}.\footnote{Není součástí výše uvedeného archivu,
    jenž neobsahuje výstupy ani intermediální soubory \pgm{*.aux},
    \pgm{*.log}, \pgm{*.toc}, \pgm{*.lof}, \pgm{*.lot}, \pgm{*.out},
    \pgm{*.bbl}, \pgm{*.blg}, \pgm{*.idx}, \pgm{*.ind}, \pgm{*.ilg}.}
\end{enumerate}

\noindent
Tato možnost je vhodná pro ty uživatele, kteří si chtějí integrovat balíček
\pkg{\ipotname} do svého \TeX{} systému, takže jej budou moci využívat pro
více projektů.

\begin{figure}[p]
\setlength{\fboxsep}{0sp}
\begin{center}
\fbox{\includegraphics[width=.86\linewidth]{MiKTeXdirs}}
\end{center}
\caption[Přidání adresáře do cest \distroMiKTeX u]%
{\label{MiKTeXdirs}Přidání adresáře do cest distribuce \distroMiKTeX .}
\end{figure}

\begin{figure}[p]
\setlength{\fboxsep}{0sp}
\begin{center}
\fbox{\includegraphics[width=.86\linewidth]{TeXLiveInstaller1}}
\end{center}
\caption[Přidání adresáře do cest \distroTeXLive]%
{\label{TeXLivedirs}Přidání adresáře do cest distribuce \distroTeXLive{} (TEXMFHOME).}
\end{figure}

Integraci provedete rozbalením prvního archivu včetně celých adresářových
stromů do libovolného adresáře,\footnote{Název a cesta opět bez znaků
  s~diakritikou a mezer.}  který následně pomocí konsole vašeho \TeX{} systému
přidáte do databáze prohledávaných cest. Kupříkladu v~konsoli distribuce
\distroMiKTeX~\cite{Sch:LIVE::MiKTeX} použijete postup znázorněný na
obrázku~\ref{MiKTeXdirs}, v~instalátoru distribuce
\distroTeXLive~\cite{TUG:LIVE::TeXLive} použijete postup znázorněný na
obrázku~\ref{TeXLivedirs}, přičemž obsah archivu byl v~obou případech rozbalen
to prázdného adresáře \pgm{texmf}. Po přidání cesty je třeba aktualizovat
databázi prohledávaných cest; příklad tohoto kroku v~\distroMiKTeX u je na
obrázku~\ref{MiKTeXrefr}, v~\distroTeXLive{} na obrázku~\ref{TeXLiverefr}.

\begin{figure}[p]
\setlength{\fboxsep}{0sp}
\begin{center}
\fbox{\includegraphics[width=.82\linewidth]{MiKTeXrefr}}
\end{center}
\caption[Aktualizace databáze cest \distroMiKTeX u]%
{\label{MiKTeXrefr}Aktualizace databáze cest distribuce \distroMiKTeX .}
\end{figure}

\begin{figure}[p]
\setlength{\fboxsep}{0sp}
\begin{center}
\fbox{\includegraphics[width=.82\linewidth]{TeXLiveShell}}
\end{center}
\caption[Aktualizace databáze cest \distroTeXLive]%
{\label{TeXLiverefr}Aktualizace databáze cest distribuce \distroTeXLive .}
\end{figure}

Alternativní možností je rozbalení prvního archivu včetně celých adresářových
stromů do některého ze systémových adresářů vaší instalace \TeX
u.\footnote{V~\distroMiKTeX u např. do \pgm{c:\textbackslash
    Users\textbackslash username\textbackslash AppData\textbackslash
    Roaming\textbackslash MiKTeX}.}  Takový adresář není třeba přidávat do
databáze prohledávaných cest (již v~ní je), opět je však třeba aktualizovat
databázi prohledávaných cest podle obrázku~\ref{MiKTeXrefr}
nebo~\ref{TeXLiverefr}.

S~obsahem druhého archivu, \pgm{\thesnamecze .zip}, pak naložíte způsobem
analogickým jako v~Sekci~\ref{installinone}. Nebude již obsahovat nepřehledné
množství systémových souborů, funkčnost však bude identická.

\section{Testování}\label{insttest}%%%%%%%%%%%%%%%%%%%%%%%%%%%%%%%%%%%%%%%%

Po instalaci libovolným z~obou výše popsaných způsobů doporučujeme otestovat
funkčnost balíčku \pkg{\ipotname} v~\TeX{} systému na vašem počítači
následovně. Aniž byste na zdrojových textech či názvech adresářů cokoli
měnili, přemístěte se do adresáře s~fiktivní závěrečnou prací \pgm{manual} a
proveďte kompilaci pomocí formátu pdf\LaTeX, poté zformátujte bibliografii
pomocí \BibTeX u, a nakonec kompilujte pdf\LaTeX em tolikrát, dokud z~logu
nezmizí hlášky typu\footnote{Jsou-li příkazy \pgm{\textbackslash cite}
  přítomny přímo v~bibliografických položkách v~\pgm{*.bib} souborech, což je
  případ tohoto manuálu, je zapotřebí dokonce dvojí zpracování \BibTeX em:
  nejprve pdf\LaTeX, poté \BibTeX, pdf\LaTeX, opět \BibTeX, a pdf\LaTeX{}
  tolikrát, dokud nezmizí uvedené hlášky.}

{\small
\begin{verbatim}
Package natbib Warning: There were undefined citations.
Package natbib Warning: Citation(s) may have changed.
(natbib)                Rerun to get citations correct.
LaTeX Warning: There were undefined references.
LaTeX Warning: Label(s) may have changed.
               Rerun to get cross-references right.
\end{verbatim}
\par}

Mějte na paměti, že některá IDE jako např. TeXworks distribuovaný jako součást
\distroMiKTeX u mají implicitně nastavenu kompilaci
pdf\LaTeX${}+{}$\MakeIndex${}+{}$\BibTeX; toto je třeba změnit na kompilaci
pouze pdf\LaTeX em a \BibTeX em podle obrázku~\ref{TeXworksCompile}: nechceme
aktualizovat soubor s~vytvořeným rejstříkem \pgm{\thesnamecze .ind} a
spouštění pomocného programu \MakeIndex{} je zbytečné. Uživatelé linuxového
shellu~\cite{Sie:1999:LinuxNutshell:} mohou kompilovat pomocí příkazového
řádku jako na obrázku~\ref{BashCompile}.

\begin{figure}[t]
\setlength{\fboxsep}{0sp}
\begin{minipage}{.48\linewidth}
\begin{center}
\fbox{\includegraphics[width=0.99\linewidth]{TeXworks_pdfLaTeX}}
\end{center}
\end{minipage}\hfill%
\begin{minipage}{.48\linewidth}
\begin{center}
\fbox{\includegraphics[width=0.99\linewidth]{TeXworks_TypesetMenu}}
\end{center}
\end{minipage}
\caption[Nastavení testovací kompilace v~IDE TeXworks]%
{\label{TeXworksCompile}Nastavení testovací kompilace v~integrovaném vývojovém
  prostředí TeXworks z~distribuce \distroMiKTeX .}
\end{figure}
\begin{figure}[t]
{\small
\begin{verbatim}
$ c=pdflatex; i=vzorova_prace; $c $i; bibtex $i; $c $i; $c $i
\end{verbatim}
\par}
\caption[Testovací kompilace v~bash]%
{\label{BashCompile}Testovací kompilace v~shellu \pkg{bash}. Počet běhů
  pdf\LaTeX u (\pgm{\$c \$i}) po sestavení bibliografie může být v~případě
  potřeby zvýšen.}
\end{figure}

Kompilace by měla proběhnout bez chyb, pochopitelně po \emph{on the fly}
doinstalaci chybějících požadovaných balíčků při první kompilaci, jak to dělá
např. distribuce \distroMiKTeX~\cite{Sch:LIVE::MiKTeX} (což může trvat nějaký
čas). Jejím výsledkem by měl být hypertextový \designation{pdf} soubor
\pgm{\thesnamecze .pdf} obsahově shodný s~tímto manuálem ve formě fiktivní
závěrečné práce.\footnote{Velikost se může lišit podle nastavené komprese
  \designation{pdf}.}

Proběhla-li kompilace bez chyb, lze bez obav začít používat balíček
\pkg{\ipotname} pro přípravu vlastní práce.

Pokud nic z~výše uvedeného nefunguje, kontaktujte lokálního \TeX ového wizarda
nebo se obraťe o~radu na \url{https://tex.stackexchange.com/} resp.
\url{https://stackoverflow.com/questions/tagged/latex}.

\section{Copyright}\label{copyright}%%%%%%%%%%%%%%%%%%%%%%%%%%%%%%%%%%%%%%%

Balíček \pkg{\ipotname} je freeware. Omezení: kýmkoli (kromě autora) provedená
jakákoli úprava nesmí být distribuována pod stejným jménem, včetně jmen
adresářů, souborů s~dokumentní třídou, bibliografickým stylem, a všech výskytů
retězce \pgm{"\ipotname "} v~nich a ve zdrojových, konfiguračních a ostatních
textových souborech knihovny a manuálu.

\endinput

%% 
%% End of file `cp_inst.tex'.
